\chapter{Objetivos de Desarrollo Sostenible}
\label{ap:ods}
% Anexo obligatorio desde las convocatorias de 2022. Reproduce la plantilla oficial de la
% ETSINF: tabla de los 17 ODS y reflexión de 500 a 1.500 palabras. El mismo texto va en
% thesis/anexo-ods.docx, el fichero suelto que se sube a Ebrón con el depósito.

\section{Grado de relación del trabajo con los ODS}
\label{sec:ods-tabla}

La Tabla~\ref{tab:ods} sitúa este trabajo frente a cada uno de los 17 Objetivos de Desarrollo Sostenible de la Agenda 2030, en los cuatro niveles que fija la plantilla de la escuela. Ningún objetivo alcanza el nivel alto, porque el trabajo no actúa directamente sobre ningún problema social ni ambiental. Tres quedan en el nivel medio, los relativos a la energía, a la innovación y al consumo responsable, y dos en el nivel bajo, la educación y la acción por el clima. La Sección~\ref{sec:ods-reflexion} justifica cada uno.

\begin{table}[htbp]
  \centering
  \caption{Relación del trabajo con los 17 ODS.}
  \label{tab:ods}
  \begin{tabular}{>{\raggedright\arraybackslash}p{8cm}cccc}
    \toprule
    Objetivos de Desarrollo Sostenible & Alto & Medio & Bajo & No procede \\
    \midrule
    1. Fin de la pobreza & & & & $\times$ \\
    2. Hambre cero & & & & $\times$ \\
    3. Salud y bienestar & & & & $\times$ \\
    4. Educación de calidad & & & $\times$ & \\
    5. Igualdad de género & & & & $\times$ \\
    6. Agua limpia y saneamiento & & & & $\times$ \\
    7. Energía asequible y no contaminante & & $\times$ & & \\
    8. Trabajo decente y crecimiento económico & & & & $\times$ \\
    9. Industria, innovación e infraestructuras & & $\times$ & & \\
    10. Reducción de las desigualdades & & & & $\times$ \\
    11. Ciudades y comunidades sostenibles & & & & $\times$ \\
    12. Producción y consumo responsables & & $\times$ & & \\
    13. Acción por el clima & & & $\times$ & \\
    14. Vida submarina & & & & $\times$ \\
    15. Vida de ecosistemas terrestres & & & & $\times$ \\
    16. Paz, justicia e instituciones sólidas & & & & $\times$ \\
    17. Alianzas para lograr objetivos & & & & $\times$ \\
    \bottomrule
  \end{tabular}
\end{table}

\section{Reflexión}
\label{sec:ods-reflexion}

Entrenar una red neuronal profunda es hoy una de las operaciones más caras del aprendizaje automático, en tiempo de máquina, en energía y en dinero. La forma habitual de saber si una configuración merece la pena sigue siendo lanzarla entera y mirar el resultado al final, y buena parte del cómputo se gasta en entrenamientos que un observador con mejor información habría detenido pronto. Este trabajo estudia una de las señales que podrían aportar esa información, las métricas calculadas sobre los gradientes de la red durante su fracción temprana, y mide si predicen la eficiencia del entrenamiento completo. Es un estudio correlacional sobre 960 entrenamientos de clasificación de imágenes, y su relación con la Agenda 2030 pasa por el cómputo que ese tipo de entrenamiento consume.

La relación más directa es con el ODS 7, energía asequible y no contaminante, y a través de él con el ODS 13, acción por el clima. Un entrenamiento que no va a converger, o que va a hacerlo más despacio de lo que su presupuesto permite, consume la misma electricidad que uno bueno hasta que alguien lo detiene. Si una medida tomada en el primer 5\,\% del presupuesto anticipara ese desenlace, el resto de la energía podría ahorrarse, y la decisión de cortar dejaría de depender de la intuición de quien mira la curva. Esa es la promesa práctica del trabajo y también su límite, porque aquí solo se mide si la señal existe, y decidir a partir de ella queda para un estudio posterior. La relación con el objetivo se mantiene con una respuesta negativa, que es la que este trabajo obtuvo. Calcular las métricas de gradiente cuesta cómputo, en esta matriz 24,1 de las 121,7 horas totales. Ninguna de ellas superó a la información gratuita de la curva de validación en la primera ventana, así que el resultado dice a quien entrena que no gaste ese cómputo. Un resultado negativo bien medido evita un consumo tan real como el que evitaría un resultado positivo. El ODS 13 se queda en el nivel bajo porque el trabajo no mide emisiones ni energía, solo horas de una GPU, y cualquier conversión a kilovatios hora sería una estimación ajena a lo medido.

El ODS 12, producción y consumo responsables, recoge la manera en que se ha gastado el cómputo del propio estudio. La matriz se diseñó una sola vez, con el tamaño mínimo que la pregunta necesitaba, y se ejecutó en una única GPU a lo largo de 121,7 horas de reloj. Antes de lanzarla se hizo un piloto de calibración por celda para fijar el presupuesto de \emph{epochs} sobre curvas de entrenamiento reales. El lanzador guarda cada entrenamiento terminado y reanuda solo lo pendiente, así que ninguno se repitió por accidente. Los resultados completos, con los dos relojes de cada entrenamiento, el de aprendizaje y el de medición, están versionados junto al código, y cualquier lector puede reanalizarlos sin volver a entrenar nada. Declarar el coste entero de la instrumentación junto a su posible beneficio es la parte del consumo responsable que un estudio como este puede cumplir por sí mismo.

El ODS 9, industria, innovación e infraestructuras, pide reforzar la investigación científica y la capacidad tecnológica. Este trabajo aporta una pieza pequeña de esa capacidad, una comparación sistemática de ocho métricas ya publicadas contra un predictor de referencia que no cuesta nada, sobre una matriz común de conjuntos de datos, arquitecturas y optimizadores. Las métricas se propusieron en artículos distintos, cada uno con su propio banco de pruebas, y hasta ahora no se habían medido todas en las mismas condiciones. El código que las calcula, con sus pruebas unitarias, queda disponible para que otro grupo lo aplique a sus propios entrenamientos o lo extienda a dominios que aquí quedan fuera, como el texto o las arquitecturas de tipo \emph{transformer}. La relación se queda en el nivel medio porque la aportación es una evaluación de herramientas existentes y no una infraestructura nueva.

El ODS 4, educación de calidad, tiene una relación indirecta. La memoria documenta el diseño, las decisiones y los errores encontrados por el camino, y el repositorio guarda los datos, el código y las pruebas que permiten repetir el estudio de principio a fin. Ese material puede servir de ejemplo trabajado de un estudio empírico reproducible para estudiantes de ciencia de datos, y en esa medida el trabajo contribuye a la formación, aunque no se haya concebido como recurso docente. Por esa razón se marca en el nivel bajo.

Los 12 objetivos restantes no proceden. El trabajo no recoge datos nuevos y trabaja con cuatro conjuntos de datos públicos de imágenes, de uso general en la investigación. No toma ninguna decisión que afecte a personas, a la salud, a la alimentación, al agua, a las ciudades, a los ecosistemas o a las instituciones. Tampoco tiene una dimensión de género ni de reducción de desigualdades más allá de la que comparte cualquier trabajo que publica su código en abierto. Marcarlos como no procedentes es más honesto que forzar una relación que la memoria no puede sostener.

La Agenda 2030 pide a la ingeniería que gaste menos para conseguir lo mismo. En el aprendizaje automático, el gasto que más crece es el del entrenamiento, y saber pronto qué entrenamientos no compensan es una de las pocas vías de ahorro que no exige cambiar el modelo ni los datos. Este trabajo mide si la señal para esa decisión existe antes de que nadie la use, y encuentra que ninguna métrica de gradiente mejora a la más barata de todas, la curva de validación.
